\section{Course Schedule I}
\label{sec:graph-course-schedule-1}

\problemstatement{Given $\textit{numCourses}$ and a list of
prerequisite pairs $[a,b]$ (meaning course $a$ requires course $b$ to
be completed first), determine whether it is possible to finish all
courses.}

\begin{patternbox}
\emph{``Can all tasks be completed given prerequisite constraints''} is
a direct restatement of Section~\ref{sec:graph-cycle-directed}: it is
possible to finish all courses if and only if the prerequisite graph
(edge $b \to a$ for every pair $[a,b]$, since $b$ must come first)
contains \textbf{no cycle} -- a cycle of prerequisites would mean some
course indirectly requires itself.
\end{patternbox}

\subsection*{Reduction}

Build a directed graph with $\textit{numCourses}$ nodes, adding an edge
$b \to a$ for every prerequisite pair $[a,b]$ (course $b$ must be
completed before $a$, so the edge points from prerequisite to
dependent). The answer is simply
$\neg\,\textit{hasCycle}(\textit{numCourses}, \textit{adj})$, using
either cycle-detection approach from Section~\ref{sec:graph-cycle-directed}
unchanged.

\subsection*{C++ implementation}

\begin{lstlisting}
#include <bits/stdc++.h>
using namespace std;

bool dfsHasCycle(int node, vector<vector<int>>& adj, vector<bool>& visited,
                  vector<bool>& pathVisited) {
    visited[node] = true;
    pathVisited[node] = true;

    for (int neighbor : adj[node]) {
        if (!visited[neighbor]) {
            if (dfsHasCycle(neighbor, adj, visited, pathVisited)) return true;
        } else if (pathVisited[neighbor]) {
            return true;
        }
    }
    pathVisited[node] = false;
    return false;
}

bool canFinish(int numCourses, vector<vector<int>>& prerequisites) {
    vector<vector<int>> adj(numCourses);
    for (auto& p : prerequisites) {
        int a = p[0], b = p[1];
        adj[b].push_back(a); // b must come before a
    }

    vector<bool> visited(numCourses, false), pathVisited(numCourses, false);
    for (int i = 0; i < numCourses; i++) {
        if (!visited[i] && dfsHasCycle(i, adj, visited, pathVisited)) {
            return false;
        }
    }
    return true;
}

int main() {
    vector<vector<int>> prerequisites = {{1, 0}, {0, 1}};
    cout << (canFinish(2, prerequisites) ? "true" : "false") << endl; // false
    return 0;
}
\end{lstlisting}

\subsection*{Dry run}

\dryrun{Take $\textit{numCourses}=2$,
$\textit{prerequisites}=[[1,0],[0,1]]$.}

Edge from pair $[1,0]$: $0 \to 1$ (course $0$ before course $1$).
Edge from pair $[0,1]$: $1 \to 0$ (course $1$ before course $0$).
This is a $2$-cycle: $0 \to 1 \to 0$. Cycle detection immediately
finds the back edge, so $\textit{canFinish}$ returns \texttt{false} --
correctly, since each course requires the other to be completed first,
an impossible circular dependency.

\begin{complexitybox}
\textbf{Time Complexity.} $O(V+E)$, inherited unchanged from cycle
detection.
\textbf{Space Complexity.} $O(V+E)$ for the adjacency list, plus
$O(V)$ for the cycle-detection arrays.
\end{complexitybox}

\begin{takeawaybox}
Recognizing that ``can all tasks be completed'' is a cycle-detection
question in disguise is the entire difficulty of this problem -- once
that reduction is seen, no new algorithm is needed at all.
\end{takeawaybox}
